% Supplementary material: moved from Section 1.2 of the main text.

\section{Reporting Dimensions and Artifact Vocabulary}
\label{supp:definitions}

\subsection{Reporting Dimensions}
\label{subsubsec:intro_def_dimensions}

A single parameter count is not enough to determine whether a model will run on a
given device. Table~\ref{tab:reporting_dimensions} lists the dimensions this
survey reports where the primary source supplies them, and states which of them
each of our claims depends on. Where a source omits a dimension, the entry is
marked as not reported rather than estimated.

\begin{table}[t]
\centering
\caption{Reporting dimensions for a small language model, the deployment question each answers, and the coverage achieved across the 209-variant inventory. ``Derived'' means the value follows arithmetically from others rather than being reported.}
\label{tab:reporting_dimensions}
\footnotesize\adjustbox{max width=\linewidth}{%
\begin{tabular}{@{}lllc@{}}
\toprule
\textbf{Dimension} & \textbf{Question it answers} & \textbf{Bounds} & \textbf{Inventory coverage} \\
\midrule
Total parameters            & How large is the download? & C2 & 209/209 \\
Active parameters per token & What does one token cost? & C1 & 36/209 stated below total; else $=$ total \\
Weight memory at a precision & Will the weights fit? & C2 & Derived \\
Peak runtime memory          & Will it fit while running? & C2 & Not reported by vendors \\
FLOPs per token             & Compute cost per token & C1 & Derived (approximate) \\
KV-cache size per token     & Cost of long context & C1, C2 & Derived from head config \\
Context length              & Largest usable input & -- & 135/209 \\
Target hardware             & What did the vendor test on? & -- & Vendor-stated only \\
\bottomrule
\end{tabular}%
}
\end{table}

\subsection{Model-Artifact Vocabulary}
\label{subsubsec:intro_def_artifacts}

Because the inventory counts artifacts rather than papers, we also fix the
following. A \emph{family} is a named line of models sharing an architecture and
a training recipe (Qwen2.5, Gemma~3). A \emph{variant} is one released
parameter size within a family. Where a vendor publishes the base and the
post-trained checkpoint of that size in one release, the inventory records one
variant and marks its post-training status as ``both'' (63 of the 209); where
only one is published, or the two were released separately, the status is
recorded as that one. The post-training status is therefore a field of the
variant, not a second count. A \emph{checkpoint}
is a specific set of weights published under a version identifier; a variant may
be revised through several checkpoints, and the inventory records the variant,
not each checkpoint. A \emph{release} is a dated publication event covering one
or more variants. A \emph{pretrained} (or base) model is the artifact produced by
pretraining alone; an \emph{instruction-tuned} model has undergone supervised
instruction fine-tuning and usually a preference-optimization stage; a
\emph{reasoning-tuned} model has additionally been trained, by distillation of
long chain-of-thought traces, reinforcement learning, or both, to emit extended
intermediate reasoning before its answer. These three post-training statuses are
not interchangeable for evaluation purposes, which is why an evaluated artifact
must be identified to the checkpoint even though the inventory counts variants, and conflating them is one of the
specific defects in the public record documented in
Section~\ref{subsec:reported_overall}.

Two conventions follow from this vocabulary and are used consistently throughout.
Family names are written as their developers write them and with a space before
the version, as in Gemma~3, Qwen2.5, Phi-4, and Llama~3.2; the original 2023
LLaMA release keeps its own capitalization, since Meta renamed the line to
``Llama'' from Llama~2 onward. Checkpoint identifiers are written in the
hyphenated form used on the model hub, as in \texttt{Gemma-3-4B-it} or
\texttt{Qwen2.5-7B-Instruct}, and are used wherever the specific artifact rather
than the family is meant.

